% 本文件由 LaTeX 排版 Agent 根据有效 translation.json 自动生成。
% author=Ambrose; work_id=3410; title=备忘录

\chapter{\WorkTitle{备忘录}}

% structure: p0001 source=h1 target=omitted reason=page-title-already-rendered-by-parent

\Person{苏马库斯}以元老院的名义向三位皇帝——\Person{瓦伦提尼安}、\Person{狄奥多西}和\Person{阿卡狄乌斯}——呈递其请愿书，但实际上仅呈递给第一位，即西方唯一皇帝。请愿书提出请求：恢复旧宗教，将\OriginalTerm{胜利祭台}{Altar of Victory}重新竖立在元老院议事厅，以便遵守古老习俗。应当效法先帝们所维护的，而非他们所废除的。国库不会受损失，而剥夺贞女和司铎的古老遗赠则是不义的，这种亵渎行为受到了诸神以饥荒的惩罚。请愿书在提出什么和隐去什么两方面都表现出高超的技巧。\par

1. 最尊贵的元老院，始终忠诚于你们，一旦得知罪行已受法律制裁，晚近的名声正被虔诚的君主净化，它便效法一个更有利时代的先例，吐露了长期压抑的悲伤，并命我再次担任代表，陈述其诉苦。然而，由于恶人之故，神圣皇帝拒绝赐予我听众；否则，正义本不会缺失，我的主上、皇帝们，大名鼎鼎的\Person{瓦伦提尼安}、\Person{狄奥多西}和\Person{阿卡狄乌斯}，常胜、常凯旋、永尊的诸位。\par

2. 因此，我身兼二职：作为你们的市政总督，我处理公共事务；作为代表，我向诸位陈明公民所托付之责。此处并无意志的冲突，因为人们如今已不再相信，若与君主意见不合便能在宫廷热忱上胜出。被爱戴、被敬畏、被尊崇，这比统治权更重要。谁能容忍私人争执损害国家呢？元老院有权谴责那些将其自身权力置于君主声誉之上的人。\par

3. 但我们的职责是为诸位的恩惠而守望。因为有什么比捍卫我们先祖的制度、国家的权利与命运更符合这个时代的荣耀呢？这荣耀尤为伟大，当你们明白，你们不可做任何违背先祖习俗之事。因此，我们要求恢复那曾长久有利于国家的宗教状况。请各教派和各意见的领袖计数：一位晚近的君主遵循其先祖的礼仪，一位更晚近的君主并未将其废除。若古宗教不足以构成先例，就让最近那位君主的默许成为先例吧。\par

4. 有谁与蛮族如此亲近，以致不需要\OriginalTerm{胜利祭台}{Altar of Victory}？今后我们将谨慎，避免此类事情的显现。但至少，那被拒绝给女神的尊荣应给予她的名号——你们永存的名声得益于胜利，并将继续得益更多。让那些从未受益于这股力量的人厌恶它吧。难道你们要拒绝抛弃那对你们的凯旋友善的庇护吗？所有人都希望拥有这股力量，没有人能否认，他所承认值得渴慕之物也当被尊崇。\par

5. 然而，即使避免此类不祥之兆的理由不够充分，至少也应当体面地避免损伤元老院议事厅的装饰。我们恳求你们，让我们这些老年人将幼时所接受之物留给后代。对习俗的热爱是强烈的。\Person{神圣君士坦提乌斯}的法令之短暂正当合理。你们应避免一切你们知道不久就被废除的先例。我们关心你们的荣耀与名声的永久性，以使将来不会发现任何需要纠正之事。\par

6. 我们在哪里宣誓遵守你们的法律和命令？用何种宗教约束才能使虚伪的心灵恐惧，不致在作证时说谎？诚然，万物充满神明，没有一处对作伪证者是安全的，但在宗教形式本身面前被催促，对产生犯罪恐惧有巨大力量。那座祭台维护全体的和谐，那座祭台诉诸每人的诚信，没有什么比我们的整个阶层仿佛在誓言的约束下颁布每一项法令更能赋予我们的法令权威。这样，对伪证将敞开大门，而此事将由我显赫的君主们决定，他们的荣誉由公共誓言捍卫。\par

7. 据说\Person{神圣君士坦提乌斯}也曾这样做过。让我们宁可效法那位君主的其他行为；若任何人此前犯过此类错误，他本不会做出任何类似之事。因为前人的跌倒纠正了后继者，改正是从对先前例子的责难中产生的。你们的祖先在如此新的事务中未能防止责备，这是可以原谅的。我们若模仿我们知道已被否定之事，同样的借口能对我们有用吗？\par

8. 陛下们，愿听一听同一位君主的其他行为，这些行为你们更值得效法吗？他未曾减少贞女的任何特权，他以贵族充任司铎职任，他不拒绝罗马礼仪的费用，他跟随欢欣的元老院穿行于永恒之城的每条街道，以平静的面容满足地观看神龛，他阅读刻在山墙上的诸神之名，询问神庙的起源，对建造者表示赞叹。虽然他本人信奉另一种宗教，但他为帝国保持了它自身的宗教，因为每人有自己的习俗，每人有自己的礼仪。神圣理智为不同的城市分配了不同的守护者和不同的崇拜。如同灵魂在婴儿出生时分别赐予，同样，命运的天才赐予各民族。在此出现了来自利益的证明，这最有力地为人向诸神作证。因为，既然我们的理性完全蒙蔽，我们更正确地获得对诸神的知识，除了从繁荣的记忆和证据，还能从哪里呢？既然漫长的岁月赋予宗教习俗以权威，我们应当忠于如此众多的世纪，追随我们的祖先，正如他们幸运地追随了他们的祖先。\par

9. 现在我们假设罗马以这样的话向你们说话：\Quote{杰出的君王，祖国的父亲，请尊重我的年岁，虔敬的礼仪使我达到这一高龄。让我沿用祖传的仪式吧，因为我不为此后悔。让我按自己的方式生活，因为我是自由的。这种敬拜曾使世界屈服于我的法律，这些神圣的礼仪曾把汉尼拔从城墙前击退，把塞农人从卡匹托尔山赶走。难道我竟为此而被存留，以致在晚年受责难吗？我将考虑何者应当整顿，但晚年的改革是迟缓而可耻的。}\par

10. 为此，我们为我们祖先和祖国的诸神祈求和平。一切敬拜理当被视为一体。我们仰望同一星辰，天空是共同的，同一世界环绕我们。各人寻求真理付出何等辛劳，又有何区别？我们无法经由一条道路达至如此伟大的奥秘；但这番讨论更适合于悠闲之人，此刻我们献上的是祈祷，而非争辩。\par

11. 削夺维斯塔贞女的特权，对您的国库有何益处？在最慷慨的皇帝治下所拒绝的，却是最俭省的皇帝所赐予的吗？她们唯一的荣誉就在于那所谓贞洁的薪俸。正如头带是她们头上的装饰，她们的尊荣也在于有闲暇从事祭献职务。她们所求的，在某种程度上不过是豁免的空名，因为她们本就贫穷而免于纳税。因此，凡削减她们生活所需的人，反而增加了她们的称誉，因为专为公益而奉献的贞洁，若没有报酬，功劳更大。\par

12. 愿这样的收益远离您国库的纯洁。好君王的收入应当增加于祭司的损失之外，而是来自敌人的战利品。有什么收益能抵偿这种憎恶？正因为贪婪的罪名不能加于你们的品性，那些古老收入被削减的人就更为不幸。因为在那些不贪图他人财物、抵制贪婪的皇帝治下，那取用者不为之所动的东西，仅仅是为了损害失去者而被取走。\par

13. 国库还保留了逝者遗嘱中留给贞女和侍奉者的土地。我恳求你们，正义的祭司，请将失去的继承权归还给你们城中神圣的人和地方。让人们无忧地订立遗嘱，知道在不贪婪的君王治下，所写下的内容将不受干扰。愿众人在这事上的幸福使你们喜悦，因为此类先例已开始扰乱临终之人。难道罗马的宗教不属于罗马法律吗？剥夺既非法律也非意外使之失效的财产，这该称作什么？被释奴领取遗赠，奴隶也不被拒绝立遗嘱的正当权利；唯独高贵的贞女和神圣礼仪的侍奉者被排除在继承财产之外。献身于贞洁、以天上的守护支持国祚、将友善的力量与你们的武器和鹰旗相连、为众人献上有效的誓愿，却无法与众人共享普通权利，这对公共安全有何益处？那么，奴隶制——对人的服役——反而是一种更好的境遇。我们损害了国家，而国家的利益从来不会是忘恩负义。\par

14. 但愿无人以为我只是在捍卫宗教事业，因为此类行为已导致了罗马民族的一切不幸。我们祖先的法律以适度的供养和公正的特权尊崇维斯塔贞女和诸神的侍奉者。这项赠予一直完好无损，直到堕落的钱商时代，他们将支持神圣贞洁的基金变成了雇佣普通脚夫的报酬。随之而来的是普遍的饥荒，欠收使各地区的希望落空。这不是大地的过错，我们不归咎于星辰的恶影响。霉病没有损害庄稼，野燕麦没有毁掉谷物；年岁因亵圣而失败，因为凡拒绝给予宗教的，也必然拒绝给予众人。\par

15. 诚然，若有任何此种灾祸的实例，让我们把这样的饥荒归咎于季节的力量。致命的风是此等荒芜的原因，人们靠树木和灌木维持生命，乡民的需求再一次转向多多纳的橡树。当公费维持宗教侍奉者时，各省遭受过怎样的类似灾祸？何时橡树为人所用而摇动？何时植物的根被拔起？何时各地的沃土同时弃绝了产出？那时，人与神圣贞女的供应是共享的。因为对祭司的供养是对大地出产的祝福，与其说是赏赐，不如说是保障。难道还有疑问吗？所赐予的是为了众人的利益，因为众人的匮乏已证明了这一点。\par

16. 但有人会说，公共支持仅仅被拒绝给予外国宗教的耗费。但愿好君王远未想到，从公共财产中给予某些人的东西竟能归于国库的权限。因为既然国家由个人组成，从国家出去的东西便重新成为个人的财产。你们统治一切，但为每个人保全他自己的东西；正义在你们心中比任意意志更有分量。请凭你们自己的慷慨判断，你们赐予他人的东西是否应当被视为公共财产。一旦为城市荣誉而给予的款项，便不再是赐予者的财产；起初是礼物，因习俗和时间而成为债务。因此，谁要是声称你们分担赐予者的责任，除非你们承担撤消礼物的憎恶，那便是在向你们灌输空洞的恐惧。\par

17. 愿所有教派的无形守护者恩待诸位殿下，尤其愿那些昔日助佑贵祖先者保卫你们，并受我们敬拜。我们恳求能延续宗教事务的此种状态，它曾为殿下们的神圣先父保存了帝国，并为那位有福的君王提供了合法嗣子。那位可敬的父亲从星空高处俯视司铎的眼泪，并视自己因他欣然遵守的惯例遭破坏而受责难。\par

18. 也请为你们的神圣兄弟改正那由他人献策之所行，掩盖那他不知道使元老院不悦的行为。因为那次使团被拒绝觐见他，唯恐舆论传至他耳。为维护前朝信誉，你们应毫不犹豫地废除那已证实非君王所为之事。\par

\SourceRecord{Translated by H. de Romestin, E. de Romestin and H.T.F. Duckworth.；Nicene and Post-Nicene Fathers, Second Series；Vol. 10.；Edited by Philip Schaff and Henry Wace.；Buffalo, NY: Christian Literature Publishing Co.,；1896.；<http://www.newadvent.org/fathers/3410.htm>.}
